\summarychapter{46}{Christianity in the West}{Decline and Reconstruction}

\begin{inbrief}
Nobody is sure where religion in the West stands. To some observers the West has left Christianity
behind; to others it is on the eve of a fresh movement of the Spirit. The chapter looks back over the
twentieth century, defines secularization, records the passage of strength from the mainline to the
conservative churches, and weighs three evangelical strategies, which it calls rescue, reclaiming and
relevance. It finds that none has reversed the decline and asks whether the Western church is dying
or taking a new shape.
\question{Is Western Christianity being defeated, or is it being rebuilt?}
\end{inbrief}

Assessments differ widely. One view holds that Christianity no longer shapes Western culture. Another
says that Evangelicals have prevailed, at least in the sense that the public now takes their beliefs,
and not liberal ones, to be what ``Christian'' means. A third admits the shrinking numbers and still
looks for a great awakening in North America. Borrowing the imagery of Revelation, the chapter says
that many church people cannot tell whether they belong to the new Israel or to the old Babylon.

Around 1900 the conservatives' opponent was liberal Christianity, which appeared to be winning the
culture. Liberals shared the age's confidence in moral and spiritual progress. They believed they
could keep the ethical cause of Jesus while leaving doctrine and liturgy aside, and some taught that
the inward experience of a Baptist and of a Buddhist is finally one. Such views, or a materialism
with no room for spirit at all, came to prevail in the universities and other leading institutions.
Conservatives replied with revival and with argument. The fundamentalists resisted as well as they
could, although they held more of modernity's assumptions than they knew. Preaching proved the
strongest reply: about a million people are thought to have answered Billy Sunday's appeals, and
Billy Graham preached on every continent. The ablest defender of the faith was a university teacher,
C.~S.~Lewis. The theologian of greatest influence was Karl Barth, who was not an Evangelical but
helped the cause when, after two world wars had shown what religion of human making can do, he
recalled church and university to revelation.

The wars discredited liberal optimism, and for a while Evangelicals seemed to have outlasted their
rivals. They then discovered other adversaries and a general cooling of living faith. Analysts
propose several explanations. An individualism absorbed in its own rights and fulfillment wears away
community and morals. Prosperity keeps raising the level of what people think they need, and rich and
poor alike pursue possessions without end; the book calls this consumerism a failure of virtue and
notes that most churches approve it and pass over what Jesus said about money. Technology has entered
daily life to an extent few foresaw. Television and the Internet belong to the last century, and the
iPhone, on sale only since 2007, is already carried by most college students, giving each of them
private and portable entertainment, information and immediate reactions from friends. It is part of
a globalization that is reshaping commerce and producing one youth culture. Its admirers compare it
to the alphabet and the printing press, which also spread knowledge. Its critics cite signs that it
shortens attention and dulls creativity, and suspect that it expresses and worsens a self-absorbed
society. Others again fault the church, which too often mirrors its culture when it should be
changing it. Secularism is one name for the whole process.

\ssection{Secularization}
The concept has a tangled history. At its core is the belief that religious matters can be neatly
separated from all others. Knowing God was once understood as the work of a lifetime, which engaged
every branch of learning and required a particular manner of life, holy practices and a new
community. It is now treated as a subject of study that calls for no conversion and no faith; some
even hold that faith makes a scholar unfit for it. Historians have traced the change to medieval
Scholasticism or to its breakdown, which left God as one cause among others, to the dry logic of
Protestant Scholasticism, and to the nation-state, which narrowed both the idea of God and the space
in which he could be served. Any of these may have contributed to making religion one specialty
beside physics and biology and a private concern.

In present usage secularism denotes the movements that seek to reduce or remove religion's
influence, so that the advance of the one is the retreat of the other. Its advocates tell a story
that comes from the Enlightenment: nature explains more each year, and in the end no one will need
the supernatural. The more aggressive of them call religion harmful as well as superfluous and blame
it for most of what has gone wrong in the world. The book grants that the church has had more
intelligent atheist opponents. But the present ones want religion gone, and would even strike
Christianity from the public record, so that it is hardly too much to call Western culture
``christophobic.''

\begin{sidebox}{Smarter Atheists}
David Bentley Hart's book \emph{Atheist Delusions}, subtitled \emph{The Christian Revolution and Its Fashionable Enemies}, argues that the angry critics of recent years get history wrong when they charge religion, and
Christianity above all, with nearly every human evil. He contrasts them with Friedrich Nietzsche
(1844--1900), who was better educated and saw further. The ``madman'' of Nietzsche's parable is not
mad because he denies God; the educated public, Nietzsche assumed, had stopped believing already. He
is treated as mad because his hearers will not face what follows from their unbelief. Nietzsche
understood that Western morality had grown from one particular idea of God and could not survive
unchanged without it. His reading of history therefore leads away from the modern critics'
conclusion: faith in God is the ground on which Western values stand, not the defect in them. Alister
McGrath suggests that the shrill tone of the new atheism betrays alarm. Christianity is expanding,
and more scientists of the first rank than ever profess belief in God, which refutes the secular
legend that education drives faith out.
\end{sidebox}

Indifference may be as serious as hostility. Society in general pays religion no attention, and
biblical phrases, once part of every educated person's speech, have dropped out of common use.
Surveys find believers less informed about basic doctrine and less attached to it than at any
earlier time. Even its defenders often justify the faith on secular grounds and make it a kind of
self-help with divine backing. Attendance and membership figures tell the same story. In Europe fewer
than one person in ten usually goes to church, in some places fewer than one in twenty, and in many
capitals the biggest congregations consist of immigrants from the Global South and are led by them.

\ssection{Shift in Western Christians and Churches}
Secularism altered the balance among the churches. The old Protestant establishment suffered most:
Episcopalians, Methodists, Presbyterians, the United Church of Christ, the larger Lutheran bodies and
the Disciples of Christ. They looked like monuments of a respectable past with nothing to say, and
from about 1965 they lost members and income; the joke was that the mainline had become the sideline.
Conservative bodies of evangelical, fundamentalist, charismatic and Pentecostal type did well on the
whole. The Assemblies of God went from 500,000 members to over four million. Roman Catholics in the
United States numbered fifty-five million by 1990, a rise due mainly to the rapid increase of the
Latino population, which would come to exceed a third of American Catholicism. The conservative
churches credited their faithfulness to Scripture. Other observers pointed to their clear answers and
their efficient organization.

\ssection{Evangelical Responses: Rescue, Reclaiming, and Relevance}
Conservative churches reacted to the secular advance in three ways. The first was to expect the end.
In a widely held scenario the world is sinking beyond recovery and has little to do with God's
purpose for the church, which will shortly be taken out of it in the Rapture. Hal Lindsey's enormous best sellers, beginning with \emph{The Late, Great Planet Earth} and continuing with titles such as \emph{Satan Is Alive and Well on Planet Earth}, taught a generation to expect Christ at any hour and to regard the repair of society as
wasted effort; the church's business was to get souls ready. Some Christian critics detected a trace
of Gnosticism in this low estimate of creation and of God's intention to redeem bodies as well as
souls. Repeated disappointments over dates drained the movement, though news from the Middle East or
about the European Union can still revive it and sell another book. Evangelicals kept the doctrine
but mostly concluded that they had applied it badly. Many of them also took part in the second
response, with some inconsistency, for in practice the three overlap.

The second response was to contest the secular trend in politics and culture. Jerry Falwell, a
Baptist pastor in Lynchburg, Virginia, launched the Moral Majority in 1979 and became the first
well-known voice of what was soon called the Religious Right, an alliance of fundamentalists,
Pentecostals, Evangelicals and Catholics who were dismayed by the country's morals and new to
political action. Allied groups included Robert Grant's Christian Voice, Beverly LaHaye's Concerned
Women for America and the Freedom Council of Pat Robertson, a television preacher who sought the
Republican nomination for president in 1988 and lost. All were convinced that secular humanism was
taking over the nation and that the press and the public schools were undermining the traditional
family.

Three events provoked them. In \emph{Roe v.~Wade} (1973) the Supreme Court ruled for an unmarried
Texas woman, named in the case as ``Jane Roe,'' that the right of privacy covers the decision to end
a pregnancy, a verdict that appalled Catholics and Protestants who believe life starts at conception.
In 1978 California voted on a measure concerning legal protection for homosexuals. Conservative
pastors who campaigned on it were told by the Internal Revenue Service that their churches might lose
tax-exempt status, and they organized the Christian Voice in order to continue. The third was the
Equal Rights Amendment. Falwell and others read the text passed by Congress as opening the way to
same-sex marriage and adoption and to conscripting women, and he led a successful fight against
ratification in Virginia. Abortion, gay rights and feminism thus brought conservative Christians
into politics, and further causes followed. Broadcasting carried the movement: Falwell's \emph{Old
Time Gospel Hour}, Robertson's Christian Broadcasting Network with its \emph{700 Club}, many smaller
stations, and \emph{Focus on the Family}, the radio program through which the child psychologist
James Dobson became a major advocate of traditional family life. The book judges that the chief
political effect was to inform and activate fundamentalists and Pentecostals, who had previously
stayed out of public affairs.

The third response avoided confrontation and tried to meet the culture halfway. In a time that
prized self-expression, some pastors designed ``user-friendly'' churches in which religion was almost
entirely a matter of personal preference. Because they wanted to evangelize, they studied their
market and shaped services that the intended audience would like and an inquirer would find
unthreatening. Sermons often took a practical, self-help line and, by a kind of allegory, turned the
Bible into a handbook of techniques for successful living.

The megachurch is the clearest example. As denominations lost public standing and faith became more
private, very large congregations attracted a rising share of worshipers, or of ``the religious
market,'' as it was bluntly put. They were free of the stale image of the denominations and they
suited popular taste. Like the seventy-six million Americans born in the years after the Second World
War, they thought of themselves as independent and individual. Their campuses were in use every day
for Bible study, support groups, children's programs, trips for the elderly and even weight-loss
classes. The chapter lists four typical features. The name avoided any denomination and used words
such as ``chapel,'' ``center'' or ``community,'' to signal a welcome for people of any background and
with any trouble, whether divorce, addiction or depression. The music was popular, energetic and
quick, and was more and more supplied by a commercial Christian recording industry. An attractive,
gifted preacher stood at the center and applied the Bible to everyday life. And everything suggested
competence and means: new buildings, professional staff, a program for each need.

Historians generally read this as one more step in the privatizing of faith. Even among thousands the
worshiper looked for help with private matters, with children, family harmony and emotional life,
and people who came from small neighborhood churches found the large ones simpler. The chapter
connects this with the American preference for private houses, cars, gardens and laundries, for
self-service shops, and, inside the home, for a room, a television, a telephone and a car for each
member of the family. Many megachurches continue to thrive, while many have stopped growing. A
younger generation may want community as a remedy for isolation. In Britain and the United States
hundreds of new congregations set out to reach postmodern young people, arguing that old methods of
ministry now conceal the gospel behind ways of thinking and speaking that no longer communicate.
These ``emerging'' churches put feeling before linear reasoning, experience before propositions,
inclusion before exclusion, and shared participation in worship before the anonymity of the
megachurch crowd.

\ssection{Troubling Signs}
Neither waiting for rescue, nor winning elections, nor adapting to the culture has turned the tide.
The evidence suggests that the ``old-time religion'' is fast becoming a minority faith. Conservatives
used to regard decline as a liberal disease. Many of their own denominations are now shrinking too,
and a pause in the losses during the early 1990s offers little reassurance. The chapter asks whether
they will go the way of the mainline, and North America the way of Europe. The Pew Research Center
reports the lowest level of religious affiliation since its records begin in 1930: a fifth of all
Americans, and a third of adults younger than thirty, claim none. The Religious Right has also lost
ground on its three issues. In June 2013 the Supreme Court struck down California's Proposition Eight
together with the federal Defense of Marriage Act, a heavy blow to legal restrictions on same-sex
marriage.

\ssection{Ongoing Struggles}
Marriage and family life have been weakened. The older convictions, that marriage comes before
living together and that it lasts, are no longer held with any consistency, even among church
members, and the book observes that the culture had given up the traditional understanding of
marriage well before the campaign for same-sex unions challenged it. One family in five now has a
single parent, and women are at least two-fifths of the labor force; both facts add to the strain on
the home.

On September 11, 2001, Americans met radical terrorism at first hand and learned that religion, in
fanatical forms too, would shape politics for years to come. Two fully fueled airliners were flown
into the twin towers of the World Trade Center in New York and destroyed them, which a bombing in
1993 had not managed to do. A third struck the Pentagon. A fourth, aimed at the White House, came
down in southwestern Pennsylvania when its passengers resisted the hijackers. President George
W.~Bush, in office for less than a year, declared war on the terrorists within days. A largely
Western coalition under his leadership entered Afghanistan, where the Taliban had sheltered and
trained them, removed the Taliban in a few months with the help of Afghan warlords and set up a new
government in Kabul. Osama bin Laden, the most dreaded of the jihadist leaders, was eventually
tracked down and shot, though not before he had inspired and directed more attacks on ``infidels.''
The coalition next invaded Iraq, said to harbor other jihadists, and overthrew Saddam Hussein. In
that bigger and stronger country the war turned out to be much harder.

\begin{sidebox}{Modernity and Islam}
Modernity is often defined in religious and political terms: a state is modern when it shelters
several religions under one national roof and stays more or less neutral among them, so that
Christians and Muslims, for example, can live together. The definition sheds light on Islam. After
the First World War nearly all Arabs found themselves in states supervised by European powers, which
they considered Christian. Some of these states were artificial and joined hostile groups under a
single flag; the arrangement was one instance of imperialism. Christians had come to terms with the
modern state long before, in the aftermath of the Reformation and the Thirty Years' War. Islam is
social and political almost by nature, and its teachers usually insist that faithfulness requires
laws enforced by government. For a period Muslims lived in modern states that Europeans had imposed
by force. Egypt's Sadat, Iran's Shah and rulers like them kept those states going by harshly
silencing radicals, and both countries passed for westernized. With independence the radicals proved
able to take power, and the Iranian revolution changed a westernized country overnight into a strict
Islamist one. Which nations will choose an Islamic state and which a modern settlement is not yet
known. For now, the sidebar says, radical Shiite voices appear to set the tone of politics and
culture, and the moderate Muslims whom most scholars believe to be the majority seem unable to make
themselves heard. Radical Islam has altered the religious map of the world and radicalized many Arab
countries. Muslim rulers have historically tolerated Christian subjects only as second-class
citizens.
\end{sidebox}

Arab Christians are paying a heavy price. They were 15 percent of the population of the Middle East
and North Africa in 1900 and are estimated at 5 percent a hundred years later, having been cowed,
driven out or killed. Of comparable importance is the arrival in the West of Muslims with a
missionary, and sometimes a militant, aim. The chapter advises vigilance toward radical Islam in both
regions and then names a worse danger, that of being unfaithful to Christ: Muslims ought to meet his
love in the concrete witness and care of Christians. Reports of conversions are difficult to verify.
There are nevertheless encouraging accounts of Muslims who encounter Jesus in dreams, of an
unexpectedly friendly hearing for some Evangelicals in Arab lands, of underground churches that
endure and of Christians who choose to live and serve among Muslims. The book counts these last among
the heroes of the missionary enterprise.

Immigration has posed another test. The Hart-Celler Act, passed with the civil rights laws, opened
the United States to people from Asia, Africa and the Middle East, and unauthorized entry from Latin
America increased the country's diversity further. Church voices had led the struggle for racial
justice, yet evangelical congregations have mostly stayed divided by race and by income, which are
linked, since a mainly white church accepts a black family more easily when its style of life is
similar. Western Christians also began to have difficulties with foreign missions just as missions
were expanding everywhere else. Budgets were cut. People asked whether humanitarian aid should take
precedence over evangelism and how much of their own culture missionaries press on others, and closer
acquaintance with other world religions was humbling. Zeal for missions had been easier when the
mission field was far away; travel and technology have now brought it next door.

\ssection{Defeated or Reconstituting?}
The chapter's answer is that reports of Christianity's death in the West may be premature. Western
churches will at least keep the influence that comes from wealth and from schools whose degrees still
carry weight in the Global South. But the climate is growing more secular and less friendly. Anyone
hoping for a restored Christendom, in which church and state cooperate as partners, will be
disappointed. Christians should expect to be pushed to the margins and to face penalties; the book
thinks prosecutions of pastors for hate speech and the loss of federal funds by Christian schools all
but inevitable. The mainline bodies will go on shrinking, some adjusting to cultural expectations
more than others, while a few may find their way back to the central tradition of the church and to
the gospel, a course known as post-liberalism. Some Pentecostal, independent and conservative
churches will keep growing. The most important gains, however, are the newly planted congregations,
which statistics badly undercount. They are flexible and mobile, gather in rented shops and in small
groups, and can scarcely be tallied.

R.~L.~Hatchett, who revised the book, adds a recollection of his own. In the modest neighborhood of
his childhood the normal churches were Baptist (his own), Methodist and Church of Christ, and their
buildings looked like churches and were finer than the houses of their members. Then a charismatic
congregation without denominational ties began to meet in a former grocery. It had a long, wordy
sign, people poorer than the rest, no connection with any seminary or mission board, and above all
the status of an oddity. Nearly fifty years later congregations of that kind are the majority, and
the near future belongs to such independent, enterprising groups.

Pentecostal and nondenominational churches are thus the new normal, and their rise marks a change in
what Western Christianity is. A time of rebuilding follows this social revolution, and most
Evangelicals will be found in these mobile churches with their many members from minorities. The
book sees something Anabaptist in them. They are not all pacifist, although many Pentecostals refused
to bear arms in the First World War, and some feel nostalgia for a Christian America; but they regard
themselves as a counterculture, concentrate on discipleship and fellowship, and cooperate to meet
human need without the political ambitions of the earlier Evangelicals. Other hopeful signs are
mentioned: a numerous body of young American Catholics who are serious about study and devotion;
student gatherings such as Passion, which brought sixty-four thousand young adults to Atlanta in 2013
and fills its conferences at home and abroad; frequent calls for renewal; and Renovar\'e, an
ecumenical fellowship devoted to spiritual formation. The chapter ends with a prayer that the Spirit
will move.

\begin{keyfacts}
\keypeople{Billy Sunday; Billy Graham; C.~S.~Lewis; Karl Barth; Friedrich Nietzsche; David Bentley
Hart; Alister McGrath; Hal Lindsey; Jerry Falwell; Pat Robertson; James Dobson; George W.~Bush; Osama
bin Laden.}
\keydates{%
  \kd{1973}{\emph{Roe v.~Wade} decision on abortion}
  \kd{1979}{Jerry Falwell founds the Moral Majority}
  \kd{1988}{Pat Robertson seeks the Republican presidential nomination}
  \kd{2001}{Terrorist attacks of September 11; war in Afghanistan}
  \kd{2007}{The iPhone introduced}
  \kd{2013}{United States Supreme Court sets aside Proposition Eight and the Defense of Marriage Act}}
\keyterms{Post-Christian; secularization; secularism; consumerism; mainline and conservative churches;
Rapture; Religious Right; Moral Majority; televangelists; megachurch; privatization of faith;
emerging churches; radical Islam; modernity; post-liberalism; new church starts.}
\end{keyfacts}
